\chapter{Introductie}
\label{chap:introductie}

Het vak Logica binnen de opleiding ICT draait om basisvaardigheden rond Computational Thinking en de onderliggende wiskundige vaardigheden. Deze vaardigheden komen in latere fases van de studie terug bij bijvoorbeeld computerorganisatie, 3D graphics en statistiek, maar zijn ook belangrijk bij het begrijpen van een computer. Daarnaast komen de vaardigheden terug in het project wat tegelijkertijd in de module draait waar logica in zit, waardoor je het geleerde direct kunt toepassen.

\subsubsection{Leerdoelen}
Dit vak heeft de volgende leerdoelen:
\begin{enumerate}
    \item Je kan berekeningen uitvoeren op het gebied van basis wiskunde, te weten: breuken, machten en factoren
    \item Je hebt kennis van lineaire vergelijkingen en lineaire stelsels en kan hiermee rekenen.
    \item Je kan goniometrie toepassen voor het berekenen van driehoeken 
    \item Je bent bekend met het verschil tussen graden en radialen en kan deze in elkaar omzetten
\end{enumerate}

\subsubsection{Wiskunde lezen}
Bij wiskundige teksten gaat het lezen veel langzamer  dan bij 'gewone' teksten. Het is niet raar als je een zin 3 keer moet lezen voor dat je hem snapt. Kortom lezen en begrijpen van wiskunde kost tijd en oefening. 

\subsubsection{Leeswijzer}
Naast theorie bevat dit document een aantal opgaven. De antwoorden vind je achter in het document. Het idee is dat de opgaven in principe genoeg oefening bieden om de stof te beheersen. 

\newpage
\subsubsection{Planning}
Onderstaand de planning voor de lessen Logica.

\begin{table}[ht]
\caption{Planning}
\renewcommand{\arraystretch}{1.4}
\setlength{\arrayrulewidth}{0.5pt}
\begin{tabular}{|c|l|c|}
\hline \rowcolor{nhl_blue} 
{\color{white}\textbf{Week}} & {\color{white}\textbf{Onderwerp}} & {\color{white}\textbf{Hoofdstuk}} \\
    \hline \rowcolor{l_blue}
    1 (36)  & --                                       &        \\ \hline 
    2 (37)  & --                                       &        \\ \hline \rowcolor{l_blue}
    3 (38)  & Haakjes wegwerken, Ontbinden in factoren & 2 en 3 \\ \hline 
    4 (39)  & Machten                                  & 4      \\ \hline \rowcolor{l_blue}
    5 (40)  & Breuken                                  & 5      \\ \hline 
    6 (41)  & Lineaire vergelijkingen en functies      & 6      \\ \hline \rowcolor{l_blue}
    \rowcolor{vakgray}
      (42)  & Vakantie                                 &        \\ \hline
    7 (43)  & 2 vergelijkingen met 2 onbekenden        & 7      \\ \hline \rowcolor{l_blue}
    8 (44)  & Goniometrie                              & 8      \\ \hline
    T1 (45) & Toets                                    &        \\ \hline \rowcolor{l_blue}
    T2b (4) & Herkansing                               &        \\ \hline
    \end{tabular}
\end{table}

% \subsubsection{Symbolen}
% In de wiskunde wordt veel gebruik gemaakt van (misschien onbekende) symbolen. Het is belangrijk om daar precies mee om te gaan. Bijvoorbeeld: $ q_0 \in Q $, wat in normale mensentaal betekent: het element $q_0$ komt uit de lijst met de naam $Q$. 
% De meeste symbolen in dit dictaat worden algemeen gebruikt in de wiskunde maar er zijn twee speciale alleen voor dit dictaat:

% \marginnote{\color{nhl_blue} \textbf{{\Large $ \Delta $ } } \ Definitie }[0cm]
% Met het teken {\color{nhl_blue} \Large $ \Delta $ }in de kantlijn wordt een definitie aangegeven.

% \marginnote{\color{nhl_blue} \textbf{{\LARGE $ \nu $ } } Voorbeeld }[0cm]
% Met het teken {\color{nhl_blue} \LARGE $ \nu $ }  in de kantlijn wordt een voorbeeld aangegeven.

% \marginnote{\color{nhl_blue} \textbf{{\LARGE $ \epsilon $ } } \ Eigenschap }[0cm]
% Met het teken {\color{nhl_blue} \LARGE $ \epsilon $ }  in de kantlijn wordt  een eigenschap aangegeven.

% In de wiskunde is het gangbaar om letters uit het Griekse alfabet te gebruiken, bijvoorbeeld $ \delta $ (delta, de Griekse d), $ \lambda $ (lambda, de Griekse l) en $\epsilon$ (epsilon, de Griekse e).

